\sequence{Arithmétique}{Arithmétique}{2--11 septembre 2026}
\section{Entiers, multiples et diviseurs}
\begin{cours}[Définitions]
$\N=\{0;1;2;\ldots\}$ est l'ensemble des entiers naturels. $\Z$ contient les entiers naturels et leurs opposés.
Pour $a,b\in\Z$, dire que $a$ est un \textbf{multiple de $b$} signifie qu'il existe $k\in\Z$ tel que $a=bk$. On dit aussi que $b$ divise $a$.
\end{cours}
$18=6\times3$ : $18$ est un multiple de $3$ et de $6$. Tout entier est multiple de $1$ et de lui-même ; $0$ est multiple de tout entier. Le seul multiple de $0$ est $0$.
\ex{Lister les diviseurs positifs de $18$ et de $24$. Écrire trois multiples de $7$, dont un négatif. $35$ est-il un multiple de $5$ ? de $6$ ? Justifier.}
\reponse[4]{Diviseurs de $18$ : $1,2,3,6,9,18$. Diviseurs de $24$ : $1,2,3,4,6,8,12,24$. Par exemple $-7,0,14$ sont des multiples de $7$. $35=5\times7$ est multiple de $5$ ; $35=6\times5+5$ n'est pas multiple de $6$.}
\begin{cours}[Somme de multiples]
La somme de deux multiples d'un entier $a$ est un multiple de $a$.
\end{cours}
\preuve[3]{Écrivons $u=ak$ et $v=a\ell$, avec $k,\ell\in\Z$. Alors $u+v=a(k+\ell)$ et $k+\ell\in\Z$. Donc $u+v$ est un multiple de $a$. Cela vaut aussi pour $a=0$.}
\rituel{Calculer $\frac23+\frac16$ ; $3(x+2)-x$ ; l'image de $2$ par $x\mapsto 2x-3$.}{$\frac56$ ; $2x+6$ ; $1$.}

\fiche{Parité et nombres premiers}
\begin{cours}[Définitions]
Un entier $n$ est \textbf{pair} s'il s'écrit $2k$, avec $k\in\Z$, et \textbf{impair} s'il s'écrit $2k+1$.
Un \textbf{nombre premier} est un entier naturel supérieur ou égal à $2$ qui possède exactement deux diviseurs positifs : $1$ et lui-même.
\end{cours}
Les nombres premiers inférieurs à $100$ sont : $2$, $3$, $5$, $7$, $11$, $13$, $17$, $19$, $23$, $29$, $31$, $37$, $41$, $43$, $47$, $53$, $59$, $61$, $67$, $71$, $73$, $79$, $83$, $89$, $97$.
\begin{cours}[Propriété à démontrer]
Le carré d'un entier impair est impair.
\end{cours}
\preuve[4]{Si $n=2k+1$, alors $n^2=(2k+1)^2=4k^2+4k+1=2(2k^2+2k)+1$. Comme $2k^2+2k\in\Z$, $n^2$ est impair. La forme finale justifie la parité ; quelques essais numériques ne suffisent pas à démontrer l'énoncé.}
\begin{methode}[Tester si un entier est premier]
Pour $n\geq2$, il suffit de tester les diviseurs premiers inférieurs ou égaux à $\sqrt n$. En effet, si $n=ab$ avec $1<a\leq b$, alors $a\leq\sqrt n$ ; $a$ possède un diviseur premier qui divise aussi $n$.
\end{methode}
\ex{Montrer que $157$ est premier et que $161$ ne l'est pas.}
\reponse[4]{Comme $12^2<157<13^2$, on teste $2,3,5,7,11$. Aucun ne divise $157$ : pour $7$, $157=7\times22+3$ ; pour $11$, $157=11\times14+3$. Donc $157$ est premier. En revanche, $161=7\times23$ est composé.}

\fiche{PGCD et fractions irréductibles}
\begin{cours}[Définition et convention]
Pour deux entiers $a,b$ \textbf{non tous deux nuls}, $\pgcd(a,b)$ est le plus grand de leurs diviseurs positifs communs. On a $\pgcd(a,b)=\pgcd(|a|,|b|)$ et $\pgcd(a,0)=|a|$ si $a\ne0$. Par convention, on pose $\pgcd(0,0)=0$ ; ce dernier cas ne relève pas de la définition par le plus grand diviseur.
\end{cours}
Par exemple, $\pgcd(18,24)=6$, et les changements de signe de $18$ ou $24$ ne modifient pas ce résultat.
\begin{cours}[Décomposition en facteurs premiers -- admise]
Tout entier naturel supérieur ou égal à $2$ est un produit de nombres premiers. Cette décomposition est unique à l'ordre des facteurs près.
\end{cours}
\ex{Décomposer $140$, $150$, $630$ et $120$ en facteurs premiers.}
\reponse[3.2]{$140=2^2\times5\times7$ ; $150=2\times3\times5^2$ ; $630=2\times3^2\times5\times7$ ; $120=2^3\times3\times5$. On divise successivement par des nombres premiers et on vérifie le produit final.}
\begin{methode}[Rendre une fraction irréductible]
Pour $b\ne0$, diviser le numérateur $a$ et le dénominateur $b$ par $\pgcd(a,b)$. La fraction obtenue a un numérateur et un dénominateur premiers entre eux.
\end{methode}
\ex{Rendre irréductibles $\dfrac{140}{630}$, $\dfrac{84}{126}$ et $\dfrac{-120}{150}$.}
\reponse[3.2]{Les PGCD valent respectivement $70$, $42$, $30$. Ainsi $140/630=2/9$, $84/126=2/3$ et $-120/150=-4/5$. Pour $140/630$, on peut aussi simplifier les facteurs $2,5,7$ dans les décompositions.}

\fiche{Vocabulaire ensembliste en contexte}
\begin{cours}[Définitions et notations]
$x\in A$ signifie que $x$ est un élément de $A$ ; $A\subseteq B$ signifie que tout élément de $A$ appartient à $B$. $\varnothing$ est l'ensemble vide.
\begin{itemize}
\item $A\cap B$ : éléments appartenant à $A$ \textbf{et} à $B$ ;
\item $A\cup B$ : éléments appartenant à $A$ \textbf{ou} à $B$ (éventuellement aux deux) ;
\item dans un ensemble de référence $E$, $\overline A=E\setminus A$ contient les éléments de $E$ qui ne sont pas dans $A$ ;
\item pour un ensemble fini, $\Card(A)$ est son nombre d'éléments.
\end{itemize}
\end{cours}
\ex{Dans $E=\{1;2;3;4;5;6\}$, on pose $A=\{2;4;6\}$ et $B=\{3;6\}$. Donner $A\cap B$, $A\cup B$, $\overline A$ et $\Card(A)$. Les assertions $2\in A$, $\{2\}\subseteq A$ et $A\subseteq B$ sont-elles vraies ?}
\reponse[3.2]{$A\cap B=\{6\}$, $A\cup B=\{2;3;4;6\}$, $\overline A=\{1;3;5\}$, $\Card(A)=3$. Les deux premières assertions sont vraies. La troisième est fausse, par exemple $2\in A$ et $2\notin B$.}
\begin{cours}[Couple et produit cartésien]
Un couple $(a,b)$ est ordonné : en général $(a,b)\ne(b,a)$. Le produit cartésien $A\times B$ est l'ensemble des couples $(a,b)$ avec $a\in A$ et $b\in B$. Pour des ensembles finis, $\Card(A\times B)=\Card(A)\Card(B)$.
\end{cours}
\ex{Écrire tous les éléments de $\{1;2\}\times\{3;4;5\}$.}
\reponse[2]{$\{(1,3);(1,4);(1,5);(2,3);(2,4);(2,5)\}$. Il y a $2\times3=6$ couples.}

\fiche{Algorithmique et réinvestissement}
\begin{methode}[Python : entier, booléen et condition]
\texttt{a \% b} est le reste de la division euclidienne pour des entiers et $b\ne0$. Un test renvoie un booléen : \texttt{True} ou \texttt{False}. L'égalité se teste par \texttt{==} ; l'affectation utilise \texttt{=}.
\end{methode}
\begin{lstlisting}
def est_multiple(a, b):
    if b == 0:
        return a == 0
    return a % b == 0
\end{lstlisting}
\ex{Prévoir les résultats de \texttt{est\_multiple(35, 5)} et \texttt{est\_multiple(35, 6)}, puis de \texttt{est\_multiple(0, 7)}, \texttt{est\_multiple(7, 0)} et \texttt{est\_multiple(0, 0)}. Expliquer le rôle de la condition.}
\reponse[3]{Les résultats sont \texttt{True}, \texttt{False}, \texttt{True}, \texttt{False}, \texttt{True}. Le cas $b=0$ évite le calcul d'un reste modulo zéro, impossible, et applique la définition : seul $0$ est multiple de $0$.}
\ex[Problème]{On dispose de $84$ cahiers et $126$ crayons. On veut former le plus grand nombre possible de lots identiques en utilisant tout le matériel. Combien de lots peut-on faire ? Que contient chacun d'eux ?}
\reponse[4]{Le nombre de lots divise $84$ et $126$. Son maximum est donc $\pgcd(84,126)=42$. Chaque lot contient $84/42=2$ cahiers et $126/42=3$ crayons. La vérification $42\times2=84$, $42\times3=126$ confirme l'utilisation de tout le matériel.}
\prof{Diagnostic : faire expliciter « divise » et « est multiple de ». Aide : commencer par tester 2, 3, 6 lots. Approfondissement : démontrer que $n(n+1)$ est pair en distinguant les deux parités de $n$.}
